\begin{pastexam}[label={px:ntu-106-dsa-02}]{臺大資工所 106 學年度 第 2 題}
（問答）

(15 points) Suppose that we have numbers between 1 and 1000 in a binary search tree and want to
search for the number 400. Determine whether the following sequences could be or could not be
the sequence of nodes examined and explain why or why not.
\begin{enumerate}[label=(\alph*)]
\item (5 points) 32 500 300 350 400
\item (5 points) 900 800 200 699 355 385 412 372 400
\item (5 points) 515 1 500 200 480 251 320 362 453 380 400
\end{enumerate}
\end{pastexam}

\begin{pastsoln}
答案：(a) 可能；(b) 不可能；(c) 可能。

\textbf{判斷的方法}。 搜尋 $400$ 時（資結書虛擬碼~9.1），每經過一個節點 $x$：$400 < x$ 就往左，
之後經過的節點都在 $x$ 的左子樹，鍵都 $< x$；$400 > x$ 就往右，之後的鍵都 $> x$。 所以走到某一步時，
下一個鍵一定落在一個開區間 $(lo, hi)$ 裡：$lo$ 是目前為止往右轉時經過的最大鍵，$hi$ 是往左轉時
經過的最小鍵。 一開始 $(lo, hi) = (-\infty, \infty)$，每經過 $x$，$400 > x$ 就把 $lo$ 改成 $x$，
$400 < x$ 就把 $hi$ 改成 $x$。 序列可能出現，若且唯若每個鍵都落在當時的區間裡，而且最後一個是
$400$。

反過來，只要每一步都落在區間裡，把序列依序插入一棵空的二元搜尋樹，每個新的鍵都會接在前一個鍵的
下面（往左或往右正好和搜尋 $400$ 的方向一致），搜尋 $400$ 走的就是這條路徑；其他的數（$1$ 到
$1000$ 裡剩下的）再插到不在這條路徑上的位置就好。 所以這個條件也是充分的。

\textbf{(a) 32 500 300 350 400}：
\begin{center}
\small
\begin{tabular}{cccc}
\toprule
經過的鍵 & 在區間裡嗎 & 方向 & 之後的區間 \\
\midrule
32 & $(-\infty, \infty)$ $\checkmark$ & $400 > 32$，往右 & $(32, \infty)$ \\
500 & $\checkmark$ & $400 < 500$，往左 & $(32, 500)$ \\
300 & $\checkmark$ & 往右 & $(300, 500)$ \\
350 & $\checkmark$ & 往右 & $(350, 500)$ \\
400 & $\checkmark$ & 找到 & \\
\bottomrule
\end{tabular}
\end{center}
每一步都在區間裡，可能。 例如依序插入 $32, 500, 300, 350, 400$ 就是這樣一條路徑：$32$ 的右子節點
$500$，$500$ 的左子節點 $300$，$300$ 的右子節點 $350$，$350$ 的右子節點 $400$。

\textbf{(b) 900 800 200 699 355 385 412 372 400}：
\begin{center}
\small
\begin{tabular}{cccc}
\toprule
經過的鍵 & 在區間裡嗎 & 方向 & 之後的區間 \\
\midrule
900 & $\checkmark$ & 往左 & $(-\infty, 900)$ \\
800 & $\checkmark$ & 往左 & $(-\infty, 800)$ \\
200 & $\checkmark$ & 往右 & $(200, 800)$ \\
699 & $\checkmark$ & 往左 & $(200, 699)$ \\
355 & $\checkmark$ & 往右 & $(355, 699)$ \\
385 & $\checkmark$ & 往右 & $(385, 699)$ \\
412 & $\checkmark$ & 往左 & $(385, 412)$ \\
372 & $\times$ & & \\
\bottomrule
\end{tabular}
\end{center}
在 $385$ 往右走了，之後的節點都在 $385$ 的右子樹，鍵必須 $> 385$；$372 < 385$，不可能出現。 所以
這個序列不可能。

\textbf{(c) 515 1 500 200 480 251 320 362 453 380 400}：
\begin{center}
\small
\begin{tabular}{cccc}
\toprule
經過的鍵 & 在區間裡嗎 & 方向 & 之後的區間 \\
\midrule
515 & $\checkmark$ & 往左 & $(-\infty, 515)$ \\
1 & $\checkmark$ & 往右 & $(1, 515)$ \\
500 & $\checkmark$ & 往左 & $(1, 500)$ \\
200 & $\checkmark$ & 往右 & $(200, 500)$ \\
480 & $\checkmark$ & 往左 & $(200, 480)$ \\
251 & $\checkmark$ & 往右 & $(251, 480)$ \\
320 & $\checkmark$ & 往右 & $(320, 480)$ \\
362 & $\checkmark$ & 往右 & $(362, 480)$ \\
453 & $\checkmark$ & 往左 & $(362, 453)$ \\
380 & $\checkmark$ & 往右 & $(380, 453)$ \\
400 & $\checkmark$ & 找到 & \\
\bottomrule
\end{tabular}
\end{center}
每一步都在區間裡，可能。 路徑左右來回很多次，看起來怪，但依序插入這 $11$ 個數得到的樹，搜尋
$400$ 正好走這條路。

這題常見的錯是只檢查相鄰兩個數的大小關係。 (b) 的 $412 \to 372$ 是往左，$372 < 412$ 沒問題；
錯的是 $372$ 和更早的 $385$ 的關係，要一路記著區間才看得出來。
\end{pastsoln}
